\section{Introduction}
Endocrine-disrupting chemicals (EDCs) are exogenous substances that interfere with hormonal signaling and are associated with various adverse health effects \citep{la2020consensus}. Their ability to disrupt sensitive hormonal balances has been linked to detrimental effects on reproductive health, endocrine system function, metabolic processes, and carcinogenic outcomes \citep{gore2015executive}. A growing body of evidence also connects specific EDCs, such as Bisphenol A (BPA), to broader population-level outcomes, including the obesity epidemic \citep{heindel2022obesity}.

Exposure to EDCs occurs through various pathways, including ingestion of pesticides via food and water, dermal contact with personal care products, and inhalation of household dust containing chemicals from furniture, electronics, and building materials (e.g., see \cite{rudel2011environmental, dodson2012endocrine}). These exposure pathways make EDCs a significant public health concern with substantial societal costs. 
Economic modeling estimates that costs related to EDC-associated diseases total approximately €163 billion in the EU and \$340 billion in the US annually in healthcare expenditures and lost productivity \citep{kassotis2020endocrine}. Regulatory efforts struggle to keep pace with scientific evidence, as even banned substances persist in legacy products and the environment \citep{sogaard2021test}. Consequently, reducing EDC-related harm cannot rely solely on top-down regulation; it also depends on consumers' perceptions of these risks and their responses in everyday choices.

Despite broad scientific consensus that many EDCs possess hazardous properties \citep{la2020consensus}, less is known about how the public translates hazard information into risk judgments across exposure pathways and national contexts. Hazard refers to a substance's capacity to cause harm, whereas risk reflects the likelihood and severity of harm under particular exposure conditions \citep{WHO2021}. Uncertainty may therefore concern hazardous properties, the relationship between exposure and effects, or the conditions under which exposure itself occurs \citep{jansen2018breaking}. Public understanding of these distinctions is incomplete. Knowledge of dose and exposure is often limited, while chemical-risk judgments are shaped by affective and intuitive cues such as perceived naturalness and processing fluency \citep{bearth2019lay,saleh2019chemophobia,song2009}. Even when people recognize that exposure affects risk, they may still treat chemicals as broadly harmful \citep{jansen2020all}, although dose-response information can alter such judgments \citep{bearth2021chemophobia}. For EDCs specifically, awareness and knowledge of exposure sources and pathways remain limited \citep{boronow2025people}.

The psychometric paradigm offers a well-established framework for analyzing public risk judgments \citep{slovic1987perception, fischhoff1978safe}. It proposes that perceived risk is structured along several attributes: voluntariness, immediacy of effects, knowledge of those exposed, scientific knowledge, controllability, newness, chronic versus catastrophic potential, dread, and severity of consequences. These attributes explain why low-probability hazards that are high in dread and uncontrollability, such as nuclear accidents, often elicit more concern than more probable but familiar hazards, such as driving \citep{slovic2012perception}. EDC-related exposure scenarios align with this framework in distinctive ways. Exposure is often involuntary or only partially controllable in everyday contexts, effects are frequently delayed and difficult to trace to specific sources, scientific knowledge continues to evolve, personal control requires deliberate effort and resources, and potential outcomes such as infertility and cancer are perceived as severe and highly dreaded. Because EDCs combine many of the psychometric features known to amplify public perceived risk, simple unidimensional measures of risk are unlikely to capture how people evaluate EDC-related risks.

Decades of risk perception research show that perceptions of risk vary systematically across individuals and social groups \citep{gustafsod1998gender}. One of the most robust findings in this literature is that gender differences in risk perception are widespread, with women typically reporting higher perceived risk than men across a range of hazards \citep{finucane2013gender}. Beyond demographic characteristics, psychological factors shape how individuals interpret information about hazards and form attitudes toward them \citep{jenkins2024integrative}. Trust in institutions responsible for managing risks -- including government agencies, scientific organizations, and industry -- is another central determinant of risk perception \citep{siegrist2000perception, siegrist2021trust}. When institutional trust is low, perceived risks tend to be inflated and official communications become less effective \citep{siegrist2008factors,slovic1993perceived}. Research focusing specifically on EDCs mirrors these general patterns, showing that public risk perceptions in this domain are shaped by a combination of demographic, cognitive, and psychosocial factors and are highly heterogeneous across populations \citep{pravednikov2024main}.

Risk perceptions are shaped by broader cultural and national contexts, as well as personal experiences \citep{weber2006experience, knuth2014risk}. Risk communication research and the Social Amplification of Risk Framework \citep{kasperson1988social,pidgeon2003social} emphasize that public responses to hazards are shaped not only by their underlying characteristics, but also by trust, institutional credibility, communication environments, and the ways in which risk signals are amplified or attenuated across social settings \citep{bearth2022social,kasperson2022social}. This is particularly relevant for EDCs, as the same substance may be embedded in different regulatory regimes, communicated through different media environments, and encountered in applications that vary in perceived controllability and familiarity. Societal values related to health, the environment, and technology, along with regulatory philosophies, media contexts, and historical experiences with environmental and technological hazards, shape how environmental risks such as EDCs are governed and communicated, contributing to systematic cross-national differences in public awareness and risk perception \citep{clahsen2019countries, renn2013cross}. Denmark, for example, operates under the European Union’s precautionary regulatory framework \citep{ecReachRestrictions} and is characterized by high institutional trust \citep{oecdTrust2024}. The United Kingdom, by contrast, now operates with greater regulatory autonomy after Brexit and adopts a different approach to chemical oversight \citep{hseUkReachExplained}. Such differences suggest that populations may vary not only in their overall concern about EDCs but also in the exposure pathways and outcomes they find most troubling, with implications for both risk communication and policy design.

Exposure pathways have not, however, been systematically examined as an organizing feature of EDC risk judgment. Chemical-risk studies have generally examined exposure within individual domains, such as food or household products \citep{bearth2019situative, jansen2020all}, while EDC-specific studies have compared chemical or product categories or assessed knowledge of possible exposure pathways \citep{pravednikov2024main,boronow2025people}, with food and water dominating the contexts examined \citep{pravednikov2024main}. It therefore remains unclear whether consumption, inhalation, and dermal exposure consistently organize risk judgments across diverse everyday scenarios and national contexts.

These gaps motivate four questions concerning the psychometric representation of EDC exposures, the organization of perceived risk by exposure pathway, individual-level correlates of EDC health beliefs, and the stability of these patterns across countries. Therefore, we ask:

RQ1: How are EDC-related exposure scenarios located within psychometric risk space?

RQ2: Do scenarios involving exposure through consumption, inhalation, and dermal application form a consistent pathway-of-exposure hierarchy?

RQ3: How are perceived severity, perceived susceptibility, and perceived overall health risk associated with sociodemographic, cognitive, and psychosocial factors?

RQ4: To what extent are these psychometric, pathway-related, and individual-level patterns stable across Denmark, the United Kingdom, and the United States?

We address these questions with a large-scale survey conducted in Denmark, the UK, and the US. We preregistered the following hypotheses.

First, we hypothesized that EDC-related risks would be perceived as involuntary, delayed, not well known to the general public, well known to scientists, relatively controllable, relatively new (except for pesticides), and associated with moderate rather than extreme perceptions of chronicity, catastrophic potential, dread, and fatality. Second, we hypothesized that activities involving exposure to EDCs through consumption (via water or food) would be perceived as riskier than activities involving exposure through skin contact or inhalation, and that dermal exposure activities would in turn be perceived as riskier than inhalation-related activities. Lastly, we hypothesized that perceived severity, perceived susceptibility, and perceived overall health risk would be associated with age, gender, education level, income level, employment status, geographical residence, perceived health, institutional trust, familiarity with EDCs, EDC knowledge, and the presence of minor children in the household.

By integrating exposure-pathway distinctions with cross-national variation in regulatory and governance contexts, the study advances understanding of how EDC-related risks are cognitively, affectively, and socially organized by the public, with direct implications for risk communication.